\begin{table}[htbp]
\centering
\small
\caption{Corpus composition and text production conditions}
\label{tab:1}
\begin{tabular}{llll}
\toprule
Corpus part & Production condition & Category & Documents \\
\midrule
machine & generation channel & real\_claude & 270 \\
machine & generation channel & deepseek\_pro & 270 \\
machine & generation channel & gpt & 270 \\
machine & generation channel & nemotron & 269 \\
machine & task mode & P1 & 360 \\
machine & task mode & P2 & 360 \\
machine & task mode & P3 & 359 \\
machine & replicate & 1 & 540 \\
machine & replicate & 2 & 539 \\
human & level 2 & Template-based materials, seo & 231 \\
human & level 0 & Free-form authorial texts, prose & 180 \\
human & level 3 & Standardised prose, science & 167 \\
human & level 1 & Journalism, news & 152 \\
human & no level assigned & Language-contact texts, translation & 73 \\
both & total & analytical corpus & 1882 \\
\bottomrule
\end{tabular}
\begin{minipage}{\linewidth}\footnotesize\vspace{2pt}
\textit{Note.} The unit of analysis is the document. Denominator: 1882 documents of the analytical corpus, 1079 machine and 803 human. Data: corpus registry, preprocessing profile prep-v5. The table reports document counts; it contains no estimates or intervals. Scale: none; the table describes composition. Abbreviations: P1, P2 and P3 are task modes, P1 — a neutral technical brief, P2 — a strict editorial brief, P3 — a task requiring the avoidance of typical machine-text features; level is the source regimentation level, assigned from the source before reading the text. Missing values: the translation stratum receives no level: for translations the regimentation is set by the source text, not by the editorial process. No correction for multiple comparisons was applied.
\end{minipage}
\end{table}
